% answers_short.tex —— 1977 年·江苏南京市初试卷 · 简短答案速查
% 由 tools/make_short_answers.py 自 answers.tex 的 \gksolhead/\gkanswer 提取，
% 并与 1977/1977_数学真题_江苏南京市初试_答案要点.md 逐题比对后生成，勿手改（改源 MD 后重跑）。

\begin{gknotice}
\textbf{使用说明}：本节只给最终答案便于快速核对，完整推导见解析版。答案为编写者独立推导，\textbf{非官方标准答案}。
\end{gknotice}

\gkpart{题目}
\begin{itemize}[leftmargin=2.6em]
  \item[\textbf{1.}] (1) $\sqrt{3}-\sqrt{2}$；(2) $1$；(3) $-3$；(4) $1$
  \item[\textbf{2.}] (1) $\log_{2}\dfrac{3.14}{\pi}<0$；(2) $\cos40^{\circ}<\sin40^{\circ}$；(3) $\left(\dfrac{2}{5}\right)^{-\frac{2}{3}}>\left(\dfrac{2}{5}\right)^{-\frac{1}{2}}$；(4) $0.\dot{3}>\dfrac{3}{10}$
  \item[\textbf{3.}] $\sin165^{\circ}=\dfrac{\sqrt{6}-\sqrt{2}}{4}$，$\cos255^{\circ}=-\dfrac{\sqrt{6}-\sqrt{2}}{4}$。
  \item[\textbf{4.}] (1) $5$；(2) $\dfrac{3}{2}$
  \item[\textbf{5.}] $x=\dfrac{3}{2}$，$y=\dfrac{1}{2}$。
  \item[\textbf{6.}] 半径为 $12\sqrt{3}\ \mathrm{cm}$。
  \item[\textbf{7.}] $45^{\circ}$。
  \item[\textbf{8.}] $x^{2}+y^{2}-4y=0$，即 $x^{2}+(y-2)^{2}=4$。
  \item[\textbf{9.}] 该船去时速度为 $12$ 公里/小时。
  \item[\textbf{10.}] (1) $\tan C=1$；(2) 锐角三角形（$\angle C=45^{\circ}$，三个内角均为锐角）；(3) 最长边为 $\dfrac{3\sqrt{5}}{5}$
\end{itemize}
